\section{Definitions and important lemmas} \label{sec:def}
\begin{defn}[graph, \cite{diestelGraphTheory2017}]
\label{defn:graph}
  Graph $G$ is a tuple $(V, E)$, where $V$ is called a vertex set and 
  $E$ are 2-element subsets of $V$.
\end{defn}

In this work we only deal with undirected graphs and for simplicity we 
shall define $(u, v) \in E$, where $u, v \in V$ as an undirected edge.

\begin{defn}[incidency, adjacency, \cite{diestelGraphTheory2017}]
    A vertex $v \in G$ is called incident to $e \in E$ if $v \in e$. We call
    two vertices $u, v \in G$ adjacent in $G$ if they are incident to the same 
    edge.
\end{defn}

In many algorithms in this paper we will need to do some pre-processing that in 
most cases will be based on reducing a graph to a simpler form, without vertices 
that have a small amount of incident edges. To make notation simpler we can introduce 
the concept of degree of a vertex.

\begin{defn}[degree, \cite{diestelGraphTheory2017}]
    We define a \textbf{degree} of a vertex $v$ in a graph $G = (V, E)$ as the number of edges 
    in $G$ that $v$ is incident to and denote it as $\deg(v)$.
\end{defn}

After reducing a graph, we get rid of some vertices and edges incident to them.
In the process of such reduction we obtain a \textit{subgraph} of the original graph.

\begin{defn}[subgraph, \cite{diestelGraphTheory2017}]
    We call graph $G' = (V', E')$ asubgraph subgraph of $G = (V, E)$ (denoted as $G' \subseteq G$) if 
    $V' \subseteq V$ and $E' \subseteq E$.
\end{defn}

When vertices are removed from a graph, all edges incident to these vertices are
removed as well. However, if we only specify the set of remaining vertices, there
is a natural way to define the resulting subgraph: we retain every edge of the original
graph whose both endpoints remain. Such a subgraph is called an \textit{induced subgraph}.
This distinction will be useful when describing the graphs obtained
during the preprocessing steps.

\begin{defn}[induced subgraph, \cite{diestelGraphTheory2017}]
Let $G = (V, E)$ be a graph and let $V' \subseteq V$. A subgraph $G' = (V', E')$ 
of $G$ is called the subgraph \textbf{induced} by $V'$ if for every pair of vertices 
$u, v \in V'$, the edge $(u, v) \in E'$ if and only if $(u, v) \in E$. 
This subgraph is denoted by $G[V']$.
\end{defn}

It is worth noting that degrees of vertices in an induced subgraph are different than in 
the original graph. If $G = (V, E)$ and $H = G[V']$ is an induced subgraph for $V' \subseteq V$, then
for $v \in H$ we can define $\deg_H(v)$ as the number of edges in $H$ that $v$ is incident to. This notation
emphasises that we are only considering the edges within the subgraph $H$, as opposed to the vertex's 
original degree in $G$.

Another concept that is connected with subgraphs is a spanning subgraph, that will be used
later when talking about spanning trees.

\begin{defn}[spanning-subgraph, \cite{diestelGraphTheory2017}]
\label{defn:spanning-tree}
Let $G = (V, E)$ be a connected graph. We call $G' \subseteq G$ a \textbf{spanning subgraph} 
    of $G$ if $V' = V$.
\end{defn}

In order to consider separators, we have to introduce the concept of connectivity 
at first.

\begin{defn}[connectivity, \cite{diestelGraphTheory2017}]
\label{defn:connectivity}
Let $G = (V, E)$ be a graph. We call $G$ \textbf{connected} if for each $u, v \in V$ there exists a path
    in $G$ connecting $u$ and $v$.
\end{defn}

As the name suggests a separator, separates the sets of vertices from each other, meaning 
there is no edge connecting any two vertices between them. After removal of a separator, 
the graph is being disconnected.
Subgraphs that are left afterwards are called connected components.

\begin{defn}[connected components, \cite{diestelGraphTheory2017}]
\label{defn:connected-component}
    Let $G = (V, E)$ be a graph. A subgraph $G' = (V', E')$ is called a connected 
    component of $G$ if $G'$ is the subgraph induced by $V'$, 
    $G'$ is connected, and $V'$ is a maximal subset of $V$ with respect to this property.
\end{defn}

After some basic graph theoretical definitions, we move on to consider different 
classes of graphs, that will appear in this paper.

One of the simplest classes of graphs are paths, which is just a sequence of 
distinct vertices connected by edges.

\begin{defn}[path, \cite{diestelGraphTheory2017}]
\label{defn:path}
    A path is a non-empty graph $P = (V, E)$ of the form $V = \{x_0, \cdots, x_k\}$,
    $E = \{(x_0, x_1), \cdots, (x_{k-1}, x_k)\}$ where the $x_i$ are distinct.
\end{defn}

If we connect the first and last vertex of a path with an edge, we create a cycle 
(complete a cycle with that edge).

\begin{defn}[cycle, \cite{diestelGraphTheory2017}]
\label{defn:cycle}
    If $P = (V, E)$ is a path of the form $V = \{x_0, \cdots, x_k\}$,
    $E = \{(x_0, x_1), \cdots, (x_{k-1}, x_k)\}$ and $k \geq 3$, then cycle is the 
    graph $C = (V, E \cup \{(x_k, x_0)\})$.
\end{defn}

After defining a cycle, we naturally may consider a graph that does not have any cycles.

\begin{defn}[forest, tree, \cite{diestelGraphTheory2017}]
\label{defn:tree}
Let $G = (V, E)$ be a graph. We call $G$ a \textbf{forest} if it does not contain any cycles (i. e. is acyclic).
We call a connected forest a \textbf{tree} $T$.
\end{defn}

We say that a tree is \textit{minimally connected}. It means that if we remove any edge, 
it will become disconnected. Also trees are \textit{maximally acyclic}, which means 
that if we add any edge to this graph, it will contain a cycle. One more important 
property of a tree is the fact, that between any two vertices, there exists one unique 
path.

Another important class of graphs that will show up when talking about a 
Vertex Cover Problem in Chapter 3 is a bipartite graph.

\begin{defn}[bipartite, \cite{diestelGraphTheory2017}]
\label{defn:bipartite}
Let $G = (V, E)$ be a graph. We call $G$ a \textbf{bipartite} graph if there exists a partition of $V$ into two subsets $X, Y$ such that
    $X \cap Y = \emptyset$ and no two vertices in the same subset are adjacent in $G$.
    To indicate that a graph is a bipartite graph we will use notation $G = B(X, Y, E)$,
    where $X, Y, E$ are defined as above.
\end{defn}


Another concept that will be useful when discussing planar graphs and their relationship
with plane graphs is that of graph isomorphism.

\begin{defn}[homomorphism, isomorphism, \cite{diestelGraphTheory2017}]
\label{defn:isomorphism}
Let $G = (V, E)$ and $G' = (V', E')$. We call a map $\phi \colon V \rightarrow V'$ a \textbf{homomorphism} if $(\phi(u), \phi(v)) \in E'$ whenever $(u, v) \in E$. If $\phi$ is a bijection and $\phi^{-1}$ is also a homomorphism, then $\phi$ is an \textbf{isomorphism}. We call $G$ and $G'$ \textbf{isomorphic} if there exists a map $\phi \colon V \rightarrow V'$ that is an isomorphism.
\end{defn}

An isomorphism provides a correspondence between the vertices of two graphs that
preserves their structure. Thus, isomorphic graphs may have different representations
or descriptions while retaining the same graph structure.

\subsection{Planar graphs}
After basic definitions for undirected graphs in general, we shall move on to definitions specific to planar graphs which are
the class of graphs that we are interested in this paper. However at first we shall start with the definition of \textit{plane}
graphs.
\\
\begin{defn}[plane graphs, \cite{diestelGraphTheory2017}]
\label{defn:plane-graph}
A \textbf{plane graph} is a pair $(V, E)$ of finite sets with the following properties:
    \begin{enumerate}
        \item $V \subseteq \mathbb{R}^2$
        \item every edge in $E$ is an arc between two vertices from $V$
        \item different edges have different sets of endpoints
        \item the interior of an edge contains no vertex and no point of any other edge
    \end{enumerate}
\end{defn}
\noindent
A plane graph defines a graph $G$ on $(V, E)$ in a natural way.
\begin{defn}[faces, \cite{diestelGraphTheory2017}]
    Let $G = (V, E)$ be a plane graph. We call the regions of $\mathbb{R}^2 \setminus G$ the \textbf{faces} of $G$.
\end{defn}
Since $G$ is bounded, exactly one face of $G$ is unbounded. This face is the outer face of $G$, while all other faces
are called inner faces. We denote a set of faces of $G$ as $F(G)$. Now we can see that because of the specific definition 
of plane graphs they have some really interesting properties that no other class has. The famous theorem for plane graph is the
Euler's formula, that connects the number of vertices, edges and faces of a plane graph in one equation.

\begin{theorem} [Euler's Formula, \cite{diestelGraphTheory2017}] \label{thm:euler}
    Let $G = (V, E)$ be a connected plane graph with faces $F$. Number of edges, vertices and faces of $G$
    is connected by the following formula:
    \begin{equation*}
        |V| - |E| + |F| = 2
    \end{equation*}
\end{theorem}

To obtain a very important lemma, which bounds the number of edges in a plane graph by the number of vertices, we have to introduce
one more lemma.

\begin{lemma} [maximal plane, \cite{diestelGraphTheory2017}] \label{lem:maximal}
    A plane graph $G = (V, E)$ with at least 3 vertices is maximally plane (it is impossible to add any more edges without violating
    planarity) when it is the plane triangulation (i.e. each face is a triangle - a cycle of length 3 on plane);
\end{lemma}

Now we can finally introduce:
\begin{corollary} [\cite{diestelGraphTheory2017}] \label{corollary:planar-bound}
    A plane graph $G = (V, E)$ with at least 3 vertices has at most $3 |V| - 6$ edges. Every plane triangulation has exactly $3 |V| - 6$
    edges.
\end{corollary}

\begin{proof}
    By Lemma~\ref{lem:maximal} we can prove only the second part of the Corollary, and the first part naturally follows (as plane triangulation
    is maximally plane, so adding more edges would destroy planarity). Let $G$ be a 
    plane triangulation. At first we shall see that if each face is a triangle, then if we can see that each face contributes exactly 3
    edges to the total number of edges. On the other hand if we look at it from the perspective of an edge, we can see that each edge is 
    incident to exactly two faces. By that double counting argument we have that $2 |E| = 3 | F |$, so $|F| = 2/3 |E|$. By putting that
    into Euler's formula we get $|V| - 1/3 |E| = 2$, which is exactly $|E| = 3 |V| - 6$.
\end{proof}

We shall see that this is a very important property of plane graphs, as they are sparse graphs, so each $O(|V| + |E|)$ algorithm on 
them is an $O(|V|)$ algorithm. Now we shall look at the concept of a minor of a graph and see a property of plane graphs that identifies
them among other classes of graphs.

\begin{defn}[edge contraction, \cite{diestelGraphTheory2017}]
\label{defn:contraction}
Let $G = (V, E)$ be a graph and let $(u,v) \in E$ be an edge. We call \textbf{edge contraction} of $(u, v)$ an operation which transforms 
    $G$ into $G'$ by merging vertices $u$ and $v$ into one vertex $v'$, such that it is adjacent in $G'$ to all the vertices that $u$ and $v$
    were adjacent to in $G$. We also remove all the loops and multi edges that are created in such operation.
\end{defn}

This edge contraction definitions leads to a definition of a minor.

\begin{defn}[minor, \cite{diestelGraphTheory2017}]
\label{defn:minor}
Let $G = (V, E)$ be a graph. A graph $H$ is a \textbf{minor} of $G$ (denoted by $H \preccurlyeq G$) if it is isomorphic to a graph $G'$ that 
    can be obtained from $G$ in a finite sequence of vertex deletions, edge deletions and edge contractions.
\end{defn}

With definitions of edge contraction and minors comes a following lemma, that will be 
important later on during discussion of Lipton Tarjan planar separator algorithm.

\begin{lemma} [\cite{liptonSeparatorTheoremPlanar1979}]  \label{lem:shrinking}
    Let $G$ be any planar graph. Shrinking any connected subgraph of $G$ to a
    single vertex preserves planarity.
\end{lemma}

By using Corollary~\ref{corollary:planar-bound} we can perform a simple test to filter out the graphs that are not plane graphs. The example of
such a graph is $K_5$ (a maximal graph on 5 vertices). It has 5 vertices and 15 edges, and by putting into formula we get $10 \leq 9$ which 
is a contradiction. Interestingly an even tighter bound than the one from \Cref{corollary:planar-bound} can be found for bipartite graphs.

\begin{corollary} \label{corollary:bipartite-bound}
    A bipartite plane graph $G = (V, E)$ with at least 4 vertices has at most $2 |V| - 4$ edges.
\end{corollary}
\begin{proof}
    In each bipartite graph each face is bounded by a cycle of length at least 4. So by counting faces, each one contributes at least 4 edges.
    On the other side each edge is incident to exactly 2 faces. So by double counting we have $2|E| \geq 4|F|$, so $|F| \leq 1/2 |E|$.
    By putting that into Euler's formula we obtain $|V| - 1/2|E| \geq 2$ and thus $|E| \leq 2|V| - 4$.
\end{proof}

Using this corollary we can show that $K_{3, 3}$ (a maximal bipartite graph with sets of sizes 3 and 3) is not a plane graph. Interestingly 
these two graphs ($K_5$ and $K_{3,3}$) are the biggest obstacles for a graph to be a plane graph. This fact is expressed by the following famous theorem.

\begin{theorem} [Wagner's theorem, \cite{wagner1937eigenschaft}] \label{thm:wagner}
    Graph $G = (V, E)$ is a plane graph if and only if its minors include neither $K_5$  nor $K_{3,3}$.
\end{theorem}

We shall see that this theorem is much stronger than what I wrote earlier. Not only are $K_5$ and $K_{3,3}$ biggest obstacles for a graph to 
be a plane graph. Actually the minors of those graphs are the only obstacles. All those properties were shown for plane graphs. Now we can 
complete everything by defining what a planar graph is.

\begin{defn}[planar embedding, \cite{diestelGraphTheory2017}]
\label{defn:embedding}
The \textbf{planar embedding} of a graph $G = (V, E)$ in a plane is an isomorphism
    of $G$ with any plane graph $H$.
\end{defn}

In later sections a very important part will be computing the embedding quickly. 
Embedding specifies a cyclic order on edges incident to each vertex, corresponding to 
their order around the vertex in a drawing. From algorithmic perspective it is 
essential, as it allows to handle geometrical properties of planar graphs programmatically.
Thus it being an essential part of the algorithm, it is important to compute it 
quickly, ideally in linear time. Luckily, there are several algorithms that do just that (see \Cref{tab:planar-embedding-algorithms}).
One of the first ones that did in that quickly was Hopcroft-Tarjan algorithm \cite{hopcroft1974efficient}, but 
the most widely used, and the one that will be used in this paper is Boyer-Myrvold 
algorithm \cite{boyer2003stop}.

\begin{table}[htbp]
\centering
\caption{Selected linear algorithms for planar embedding.}
\label{tab:planar-embedding-algorithms}
\begin{tabularx}{\textwidth}{llX}
\hline
\textbf{Authors} & \textbf{Year} & \textbf{Original work} \\
\hline
Hopcroft and Tarjan
& 1974
&  \cite{hopcroft1974efficient}
\\
Booth and Lueker
& 1976
& \cite{booth1976testing}
\\
Chiba, Nishizeki, Abe, and Ozawa
& 1985
&  \cite{chiba1985linear}
\\
Boyer and Myrvold
& 1999
& \cite{boyer2003stop}
\\
\hline
\end{tabularx}
\end{table}

\begin{defn} [Planar Graph] \label{defn:planar}
    Graph $G = (V, E)$ is called \textbf{planar} if there exists a planar embedding that maps it to a plane graph $H$.
\end{defn}

\subsection{Hard problems and planar graphs}
From the previous section we can see that planar graphs are a very special class of graphs. Some of their properties actually help reduce 
the time complexity for some problems by a large margin. The example is maximum-clique problem which asks to find the biggest clique in 
a graph. For any abstract graph, that problem is $\mathbb{N}\mathbb{P}$-complete.
However for planar graphs it can be done in $O(|V|)$ time. It is due to the fact 
that the maximum size of a clique in planar graphs is 4, which comes directly from Theorem~\ref{thm:wagner}. 
The algorithm for that problem is based on the following lemma.
\begin{lemma} \label{lem:small-deg}
    Let $G = (V, E)$ be a planar graph. There exists a vertex $v$ such that $\deg(v) \leq 5$.
\end{lemma}
\begin{proof}
    Let's suppose that the minimum degree of $G$ is 6. Then we can once again construct a double counting argument:
    each edge is incident to exactly two edges, and each vertex is incident to at least 6 edges. So we have 
    $2|E| \geq 6|V|$, then $|E| \geq 3|V|$. Now we can use Corollary~\ref{corollary:planar-bound} and conclude that
    $3|V| \leq 3|V| - 6$ which leads to contradiction.
\end{proof}

Another application of this lemma is Graph Coloring problem which is $\mathbb{N}\mathbb{P}$-complete in general
case.
\begin{defn}[vertex coloring, \cite{diestelGraphTheory2017}]
    A vertex coloring of a graph $G = (V, E)$ is a map $c \colon V \rightarrow S$ such that $c(v) \neq c(w)$
    whenever $v$ and $w$ are adjacent. The element of the set $S$ are called colors. Graph Coloring problem 
    is finding the minimum size $S$ such that $c$ is a correct coloring.
\end{defn}

Because of \Cref{lem:small-deg} we can easily prove that for planar graphs it is sufficient to use 6 colors. However, that is 
not the smallest number possible. In 1977 it was proven by Haken and Appel \cite{appel1977solution} that 
each planar graphs is $4$-colourable, moreover it can be done in polynomial time. 
And although most of the $\mathbb{N}\mathbb{P}$-complete problems are still in this complexity class for planar graphs 
(for example determining whether a planar graph is $3$-colourable is still $\mathbb{N}\mathbb{P}$-complete \cite{gareyJohnsonStockmeyer1976}), there is another opportunity
to reduce the time of computation. We can resort to approximation algorithms that will utilize planar separator theorem, that is the 
main topic of this paper.

\subsection{Approximation algorithms}
When we can exchange the optimality of the solution for the problem to gain some computation time, we shall take into 
consideration approximation algorithms. The idea of such an algorithm is that the found solution will be worse than the optimal 
by some factor, in exchange for better time complexity. In some cases one can find a quicker algorithm for a polynomial problem, but 
sometimes we can find a polynomial approximation algorithm to an $\mathbb{N}\mathbb{P}$-complete problem which is a great improvement in terms of time 
complexity.

\begin{defn} \cite{williamson2011design} \label{defn:opt-problem}
    An $\alpha$-approximation algorithm for an optimization problem is a polynomial time algorithm that for all instances of a problem 
    produces a solution whose value is within a factor of $\alpha$ of the value of an optimal solution.
\end{defn}

For an $\alpha$-approximation algorithm, we shall call $\alpha$ the approximation ratio, and $1 - \alpha$ the approximation error.
For example, the $4$-coloring algorithm for planar graphs described in the previous section yields a $4/3$-approximation algorithm for the minimum vertex coloring problem.

We can also look for something better than a constant approximation ratio. There are some problems for which we can get a solution 
arbitrarily close to the optimal solution. These problems have polynomial-time approximation schemes.

\begin{defn}
    A polynomial-time approximation scheme is a family of algorithms $\{A_{\epsilon}\}$, where there is an algorithm for each $\epsilon > 0$,
    such that $A_{\epsilon}$ is a $(1 + \epsilon)$-algorithm (for minimization problems) or $(1 - \epsilon)$-algorithm (for maximization problems).
\end{defn}

In Chapter 3 we shall explore how planar separators can be utilized to create approximation schemes for two problems: an already polynomial-time 
problem of finding maximum matching for which we shall find an $O(n \log n)$ algorithm, and $\mathbb{N}\mathbb{P}$-complete problem of finding a minimum vertex cover.
